We agree that the city-pair observations cannot be assumed to be strictly independent and identically distributed, because each city occurs in multiple dyads. We therefore quantified three sources of dependence rather than applying conventional i.i.d.-based tests. At $K=1{,}000$, CI decreased with geographic distance ($r_s=-0.179$, one-sided city-label permutation $P<0.001$), and same-country pairs had higher CI than cross-country pairs ($\Delta\mathrm{CI}=0.103$, $P_{\mathrm{perm}}<0.001$). Moreover, source- and target-city fixed effects jointly explained $46.4\%$ of the variation in directed CI, confirming systematic dependence among observations that share a city. We then fitted multiple linear regressions with log geographic distance and same-country status as covariates:

\begin{equation}
CI_{ij}=\beta_0+\beta_1 HistoricalConnection_{ij}+\beta_2\log(Distance_{ij})+\beta_3 SameCountry_{ij}+\varepsilon_{ij}.
\end{equation}

Rather than using conventional $P$ values that assume independent observations, significance was assessed using 9,999 simultaneous row--column city-label permutations. The association between a shared historical regime and CI remained significant ($\beta=0.069$, one-sided $P_{\mathrm{perm}}=0.012$), whereas the direct-tie coefficient remained positive but was not significant after adjustment ($\beta=0.034$, $P_{\mathrm{perm}}=0.173$). These results confirm that strict i.i.d. inference is inappropriate for the dyadic data and motivate the use of city-label permutation tests and explicit geographic and national controls.

\begin{figure}[!htbp]
    \centering
    \includegraphics[width=\linewidth]{Figures/Figx_dyadic_dependence_diagnostics_k1000.pdf}
    \caption{\textbf{Dependence diagnostics and controlled historical associations at $K=1{,}000$.} \textbf{a,} Covered Index across geographic-distance groups, with Spearman correlation assessed using a city-label permutation test. \textbf{b,} Covered Index for cross-country and same-country city pairs. \textbf{c,} Variance in directed CI explained by source-city, target-city, and joint source--target fixed effects. \textbf{d,} Standardized coefficients for direct ties and shared historical regimes from separate multiple linear regressions controlling for log geographic distance and same-country status. Boxes show the interquartile range, white lines indicate medians, and whiskers extend from the 5th to the 95th percentile. All permutation $P$ values are one-sided and based on 9,999 simultaneous city-label permutations.}
    \label{fig:dyadic-dependence-diagnostics}
\end{figure}
